\section{Introdução}

O Zovic é uma plataforma de streaming de música desenvolvida na disciplina de
Engenharia de Software. Visitantes navegam pelo catálogo, assinantes ouvem
músicas, montam playlists, comentam e avaliam faixas, e um gestor pode dividir
um plano familiar com até quatro pessoas. Um administrador mantém o catálogo e
acompanha as estatísticas de uso.

Todos esses dados ficam num único banco PostgreSQL. Este relatório descreve
como ele foi modelado e como é mantido:

\begin{itemize}
  \item o \textbf{modelo conceitual} (\cref{sec:conceitual}), que mostra as
    entidades do domínio e como elas se relacionam, sem detalhes de
    implementação;
  \item o \textbf{modelo relacional} (\cref{sec:relacional}), que mapeia o
    conceitual em tabelas, chaves primárias e chaves estrangeiras;
  \item as \textbf{restrições de integridade} (\cref{sec:integridade}) que o
    próprio banco garante;
  \item a \textbf{evolução do esquema} por migrations (\cref{sec:evolucao}) e
    o \textbf{acesso aos dados} pela API (\cref{sec:acesso});
  \item as \textbf{limitações conhecidas} e o que ainda está planejado
    (\cref{sec:limitacoes}).
\end{itemize}

A fonte da verdade do esquema são as migrations em \campo{db/migrations/}, no
repositório do projeto. Este relatório acompanha a última delas,
\campo{V261117\_170000\_\_Adiciona\_imagem\_da\_playlist\_e\_do\_artista.sql}.
